\documentclass[11pt,a4paper]{article}

% ── Packages ──────────────────────────────────────────────────────────────
\usepackage[utf8]{inputenc}
\usepackage[T1]{fontenc}
\usepackage{lmodern}
\usepackage[margin=2.5cm]{geometry}
\usepackage{parskip}
\usepackage{enumitem}
\usepackage{booktabs}
\usepackage{array}
\usepackage{graphicx}
\usepackage{xcolor}
\usepackage{tcolorbox}
\usepackage{hyperref}
\usepackage{fancyhdr}
\usepackage{titlesec}

% ── Colors (Okabe–Ito colorblind-safe palette) ───────────────────────────
\definecolor{cOrange}{HTML}{E69F00}
\definecolor{cBlue}{HTML}{0072B2}
\definecolor{cSky}{HTML}{56B4E9}
\definecolor{cVermillion}{HTML}{D55E00}
\definecolor{cPurple}{HTML}{CC79A7}
\definecolor{cGray}{HTML}{999999}

% ── tcolorbox styles ─────────────────────────────────────────────────────
\tcbuselibrary{skins,breakable}

\newtcolorbox{keyinsight}[1][]{%
  enhanced, breakable,
  colback=cSky!8, colframe=cBlue,
  fonttitle=\bfseries, title=Key Insight,
  left=4pt, right=4pt, top=4pt, bottom=4pt, #1}

\newtcolorbox{managerbox}[1][]{%
  enhanced, breakable,
  colback=cOrange!8, colframe=cOrange,
  fonttitle=\bfseries, title=Manager's Briefing,
  left=4pt, right=4pt, top=4pt, bottom=4pt, #1}

\newtcolorbox{pitfall}[1][]{%
  enhanced, breakable,
  colback=cVermillion!8, colframe=cVermillion,
  fonttitle=\bfseries, title=Common Pitfall,
  left=4pt, right=4pt, top=4pt, bottom=4pt, #1}

% ── Header / Footer ──────────────────────────────────────────────────────
\pagestyle{fancy}
\fancyhf{}
\lhead{\small Stable Diffusion from Scratch}
\rhead{\small A Decision-Maker's Guide}
\cfoot{\thepage}
\renewcommand{\headrulewidth}{0.4pt}

% ── Section styling ──────────────────────────────────────────────────────
\titleformat{\section}{\Large\bfseries\color{cBlue}}{\thesection}{1em}{}
\titleformat{\subsection}{\large\bfseries\color{cBlue!80!black}}{\thesubsection}{1em}{}
\titleformat{\subsubsection}{\normalsize\bfseries\color{cBlue!60!black}}{\thesubsubsection}{1em}{}

\hypersetup{
  colorlinks=true,
  linkcolor=cBlue,
  urlcolor=cBlue!80!black
}

% ══════════════════════════════════════════════════════════════════════════
\begin{document}

\begin{titlepage}
\centering
\vspace*{3cm}
{\Huge\bfseries Stable Diffusion from Scratch\\[0.4em]}
{\LARGE\color{cBlue} A Decision-Maker's Guide}

\vspace{2cm}
{\Large Lecture Notes}

\vspace{1.5cm}
{\large CAS BMAI --- ETH Z\"urich}

\vspace{3cm}
\begin{tabular}{rl}
\textbf{Format:} & 5 interactive sessions, approx.\ 2.5 hours total \\
\textbf{Prerequisites:} & None (ability to click ``Run'' in Google Colab) \\
\textbf{Audience:} & Managers and non-technical leaders \\
\end{tabular}

\vfill
{\small \today}
\end{titlepage}

\tableofcontents
\newpage

% ══════════════════════════════════════════════════════════════════════════
\section{Introduction}
\label{sec:intro}

Text-to-image generation---the ability to produce photorealistic pictures
from a short textual description---has moved from a research curiosity to a
strategic capability in less than three years.  Products such as DALL-E,
Midjourney, and Stable Diffusion now power content pipelines in marketing,
e-commerce, entertainment, and product design.

These lecture notes accompany a hands-on course in which participants
\emph{build} a miniature text-to-image system themselves, step by step.
The system is small enough to train on a laptop yet architecturally
faithful to the full-scale models used in production.  The goal is not to
turn managers into machine-learning engineers but to equip them with a
\emph{working mental model} of the technology---one precise enough to
evaluate vendors, diagnose failures, and ask engineering teams the right
questions.

\subsection{The Co-Pilot Metaphor}

Throughout the course we use a single, consistent analogy that maps every
technical component to a familiar concept:

\begin{center}
\renewcommand{\arraystretch}{1.3}
\begin{tabular}{>{\bfseries}l l}
\toprule
Technical Component & Metaphor \\
\midrule
Diffusion model (U-Net) & The \emph{driver}---knows how to navigate but not where to go \\
Text prompt & The \emph{destination} \\
CLIP text encoder & The \emph{co-pilot} reading the map \\
Cross-attention & The driver asking the co-pilot at every turn \\
Guidance scale & How \emph{loudly} the co-pilot speaks \\
Negative prompt & The co-pilot saying ``NOT that way'' \\
Random seed & The random road network \\
\bottomrule
\end{tabular}
\end{center}

By the end of Session~5 you will have assembled all parts of this
``car'' and driven it yourself.

\subsection{Toy Dataset}

All experiments use a \textbf{programmatically generated} dataset of
coloured geometric shapes on a 32$\times$32 pixel canvas:

\begin{itemize}[nosep]
  \item \textbf{Shapes:} circle, square, triangle
  \item \textbf{Colours:} red, blue, green
  \item \textbf{Extensions (Session~5):} small/large, centred/offset
\end{itemize}

Because every image is generated with known ground truth, we can measure
model quality objectively at every stage---an advantage rarely available
with real-world data.


% ══════════════════════════════════════════════════════════════════════════
\section{Session 1: The Raw Ingredients---Images and Text as Numbers}
\label{sec:session1}

\textit{Approx.\ 25 minutes.  Metaphor addition: ``Here are the parts of the car and the map.''}

\subsection{Images as Tensors}

A digital image is stored as a three-dimensional grid of numbers---a
\textbf{tensor} of shape $H \times W \times C$, where $H$ is the height
in pixels, $W$ the width, and $C$ the number of colour channels (three
for RGB images: red, green, blue).  Each entry is a value between 0 and 1
representing the intensity of one colour at one pixel location.

For the toy dataset, every image is a $32 \times 32 \times 3$ tensor
containing $32 \times 32 \times 3 = 3{,}072$ numbers.

\subsection{Text as Embedding Vectors}

Computers cannot process words directly.  Instead, every word (or
sub-word token) is mapped to a list of numbers called an
\textbf{embedding vector}.  Words with similar meanings end up close
together in this numerical space; unrelated words end up far apart.

The notion of ``closeness'' is quantified by \textbf{cosine similarity}:
two vectors that point in roughly the same direction have a cosine
similarity near~1; orthogonal vectors score~0; opposite vectors score~$-1$.

\subsection{The Gap}

At this stage, image tensors and text embeddings live in
\emph{completely different} numerical spaces.  A raw pixel vector and
a word embedding cannot be compared.  Bridging this gap is the central
problem that the next session addresses.

\begin{managerbox}
\textbf{Why this matters.}  When an AI vendor says their model
``understands'' images and text together, they are claiming to have
bridged exactly this gap.  The quality of that bridge---how well visual
and textual representations are aligned---determines whether the model
produces what you asked for or something subtly (or wildly) wrong.
\end{managerbox}

\subsubsection*{Real-World Story: The \$180K Print Disaster}

A marketing team approved AI-generated product images on screen and sent
them directly to a print shop.  The images were encoded in sRGB (the
standard for screens); the printer expected CMYK.  Colour shifts in the
converted files ruined an entire print run.  \emph{Lesson:} image
representations (colour spaces, resolutions, file formats) are not
interchangeable, and automated pipelines must account for conversions.


% ══════════════════════════════════════════════════════════════════════════
\section{Session 2: The Translator---CLIP}
\label{sec:session2}

\textit{Approx.\ 30 minutes.  Metaphor addition: ``The co-pilot learns to read the map.''}

\subsection{The Core Idea: Contrastive Learning}

\textbf{CLIP} (Contrastive Language--Image Pre-training) is a model that
learns to place images and their corresponding text descriptions
\emph{near each other} in a shared numerical space, while pushing
mismatched pairs apart.

Training proceeds as follows:

\begin{enumerate}[nosep]
  \item Collect a batch of image--text pairs.
  \item Pass each image through an \textbf{image encoder} to obtain an
        image embedding.
  \item Pass each text through a \textbf{text encoder} to obtain a text
        embedding.
  \item Compute the cosine similarity between every image embedding and
        every text embedding in the batch.
  \item Adjust the encoders so that \emph{matching} pairs score high and
        \emph{non-matching} pairs score low.
\end{enumerate}

The loss function that drives this adjustment is called \textbf{InfoNCE}.
Intuitively, it rewards the model for ranking the correct match above all
distractors.

\begin{keyinsight}
After training, CLIP provides a \emph{common language} for images and
text.  You can compare an image to a sentence---or two images, or two
sentences---simply by measuring the distance between their embedding
vectors.
\end{keyinsight}

\subsection{Building a Mini-CLIP}

In the hands-on exercise, participants train a small CLIP model
($\sim$200K parameters) on the toy dataset.  The text encoder converts
a 12-token vocabulary into embedding vectors; the image encoder converts
flattened pixel arrays through a short sequence of transformations.
Training takes approximately 15 seconds on a CPU.

\subsection{Productive Failure: Memorisation vs.\ Generalisation}

When Mini-CLIP is trained on only three image--text pairs it memorises
them perfectly---but fails on any new example.  This demonstrates a
fundamental tension in machine learning: a model must see enough variety
during training to learn \emph{general} patterns rather than specific
examples.

\begin{managerbox}
\textbf{Vendor question.}  ``How large and diverse is your training
dataset?''  A model trained on narrow data will excel on benchmarks that
resemble its training set and fail unpredictably elsewhere.
\end{managerbox}

\subsubsection*{Real-World Story: The CLIP Bias Problem}

Researchers discovered that CLIP embeddings encode societal biases present
in their web-scraped training data.  For example, certain professions were
more closely associated with specific genders or ethnicities in embedding
space.  \emph{Lesson:} the bridge between text and images is only as fair
as the data used to build it.  Bias audits are essential before
deployment.


% ══════════════════════════════════════════════════════════════════════════
\section{Session 3: The Engine---Diffusion}
\label{sec:session3}

\textit{Approx.\ 35 minutes.  Metaphor addition: ``The driver learns to navigate by being dropped at random locations.''}

\subsection{Forward Process: Adding Noise}

The diffusion framework rests on a surprisingly simple idea.  Start with
a clean image and gradually add random noise, step by step, until the
image becomes pure static.  This is the \textbf{forward process}---it is
easy, deterministic given the noise schedule, and requires no learning.

At each step~$t$, the image becomes a little noisier.  By the final step,
no trace of the original content remains.

\subsection{Reverse Process: Denoising}

The creative act is running this process \emph{in reverse}: starting from
pure noise and removing noise step by step until a clean image emerges.
A neural network---the \textbf{denoiser}---is trained to predict the noise
that was added at each step.  Given a noisy image and the current step
number, it outputs an estimate of the noise, which is then subtracted.

Repeating this subtraction across all steps (typically 50 in the demo)
transforms random static into a coherent image.

\subsection{The Noise Schedule}

The rate at which noise is added (and, conversely, removed) is controlled
by a \textbf{noise schedule}.  The course uses a \textbf{cosine schedule},
which adds noise slowly at first, accelerates in the middle, and tapers
at the end.  The schedule strongly influences image quality: too
aggressive and fine details are destroyed before the model can learn them;
too gentle and the model wastes capacity on near-clean images.

\subsection{Productive Failure: Why Architecture Matters}

Participants first attempt denoising with a na\"ive fully connected
network (an MLP).  This network can learn colours but produces blobby,
shapeless outputs.  The reason: an MLP treats every pixel independently
and has no notion of \emph{spatial neighbourhood}.

The fix is to switch to a \textbf{U-Net}---a convolutional architecture
that processes local patches and preserves spatial structure through
skip connections.  The U-Net used in the demo has approximately 150K
parameters and trains in under three minutes.

\begin{keyinsight}
The architecture of the denoiser determines \emph{what kinds of
structure} the model can learn.  Choosing the right architecture is not
a minor implementation detail---it is a strategic decision that shapes
output quality.
\end{keyinsight}

\begin{managerbox}
\textbf{Vendor question.}  ``What is the architecture of your denoiser,
and what are its known limitations?''  A model that generates beautiful
landscapes may fail at fine text, small faces, or hands---often because
of architectural choices in the denoiser.
\end{managerbox}

\subsubsection*{Real-World Story: The Deepfake Dilemma}

Diffusion-based generators leave subtle fingerprints in the
\emph{frequency domain} of their outputs---patterns invisible to the
human eye but detectable by specialised analysis tools.  Forensic teams
exploit these artefacts to distinguish AI-generated images from
photographs.  \emph{Lesson:} every generative architecture has a
signature; understanding that signature is essential for both quality
assurance and content authenticity.


% ══════════════════════════════════════════════════════════════════════════
\section{Session 4: The Conversation---Cross-Attention}
\label{sec:session4}

\textit{Approx.\ 30 minutes.  Metaphor addition: ``At every turn, the driver asks the co-pilot.''}

\subsection{From Unconditional to Conditional Generation}

The diffusion model from Session~3 generates images, but it has no way
to steer \emph{what} it generates.  The driver can navigate but does not
know the destination.  We need a mechanism for the text prompt (via CLIP
embeddings) to influence the denoising process at every step.

\subsection{How Cross-Attention Works}

Cross-attention is the mechanism by which image features ``consult'' text
features during denoising.  At each layer of the U-Net where
cross-attention is applied:

\begin{enumerate}[nosep]
  \item The image features produce a set of \textbf{queries} (Q):
        ``What should I look like here?''
  \item The text features produce \textbf{keys} (K) and
        \textbf{values} (V): ``Here is what the prompt says.''
  \item Each query is compared to all keys via a similarity score.
  \item The resulting \textbf{attention weights} determine how much each
        text token influences each spatial location in the image.
  \item The weighted combination of values is injected back into the
        image features.
\end{enumerate}

This process repeats at every denoising step, allowing the text prompt to
continuously guide image generation.

\begin{keyinsight}
Cross-attention is not a one-time instruction; it is a
\emph{continuous conversation} between the image being generated and the
text prompt.  At every denoising step and at every spatial location, the
image ``asks'' the text what to do.
\end{keyinsight}

\subsection{Visualising Attention}

Attention weights can be visualised as heatmaps showing which parts of
the image attend to which words in the prompt.  In the demo, participants
observe that the word ``red'' activates spatial locations where the shape
is being drawn, while ``circle'' activates the shape's boundary.  This
transparency is a powerful diagnostic tool.

\subsection{Productive Failure: Uniform Attention}

When attention weights are forced to be uniform (every text token
influences every pixel equally), the output collapses into a blurry,
colour-averaged blob.  Spatial selectivity---the model's ability to
route different words to different image regions---is essential for
coherent generation.

\begin{managerbox}
\textbf{Vendor question.}  ``Can you show me attention maps for a
failed generation?''  Attention visualisation reveals \emph{why} a
model produced a wrong output, not just \emph{that} it did.  This is
invaluable for debugging and quality assurance.
\end{managerbox}

\subsubsection*{Real-World Story: The Compositional Failure}

An e-commerce company used an AI model to generate product images showing
``a leather bag on a wooden table.''  The model sometimes swapped
materials---producing a wooden bag on a leather table.  This
\emph{attribute binding} failure is a known limitation of
cross-attention: when multiple objects and attributes are present, the
attention mechanism can mis-route attributes to the wrong objects.


% ══════════════════════════════════════════════════════════════════════════
\section{Session 5: The Road Trip---The Full Pipeline}
\label{sec:session5}

\textit{Approx.\ 30 minutes.  Metaphor: ``Start the engine, set the destination, drive.''}

\subsection{Assembling the Pipeline}

In this final session, all components are wired together into a complete
text-to-image system:

\begin{enumerate}[nosep]
  \item The \textbf{text prompt} is encoded by the trained Mini-CLIP
        text encoder into an embedding vector.
  \item A \textbf{random noise image} is sampled (determined by the
        seed).
  \item Over 50 denoising steps, the \textbf{U-Net} progressively
        removes noise, consulting the text embedding via
        \textbf{cross-attention} at each step.
  \item The final output is a clean, generated image matching the prompt.
\end{enumerate}

The climactic moment of the course: participants type a prompt such as
``red circle,'' and watch the system they built transform pure noise into
a recognisable red circle over a 50-step filmstrip.

\subsection{Classifier-Free Guidance}

In practice, the denoiser runs twice at each step: once \emph{with} the
text prompt and once \emph{without}.  The difference between the two
outputs indicates the ``direction'' imposed by the text.  This difference
is then amplified by a factor called the \textbf{guidance scale}.

\begin{itemize}[nosep]
  \item \textbf{Low guidance} ($\sim$1): the model largely ignores the
        prompt, producing diverse but potentially off-topic images.
  \item \textbf{High guidance} ($\sim$7--15): the model follows the prompt
        closely, producing sharp and on-topic results but with less
        diversity.
  \item \textbf{Very high guidance} ($>$20): outputs become oversaturated
        and artefact-prone.
\end{itemize}

In the co-pilot metaphor, the guidance scale controls how loudly the
co-pilot speaks.  Too quiet and the driver wanders; too loud and the
driver oversteers.

\subsection{Seed Control and Reproducibility}

The \textbf{random seed} determines the initial noise image.  Fixing the
seed while varying the prompt allows controlled comparison of how
different text inputs steer the same starting point.  Conversely, fixing
the prompt while varying the seed reveals the range of outputs the model
can produce for identical instructions.

\begin{managerbox}
\textbf{Why seeds matter for business.}  Reproducibility is essential
for auditing, debugging, and regulatory compliance.  If a generated
image causes a problem, you need to be able to recreate it
deterministically.  Always require that your AI pipeline logs the seed
alongside every generated asset.
\end{managerbox}

\subsection{Productive Failure: Compositional Complexity}

When the dataset is extended from 9 classes (3 shapes $\times$ 3 colours)
to 36 classes (adding size and position), the model begins to struggle
with novel combinations it has seen rarely during training.  This
demonstrates the fundamental \textbf{combinatorial challenge} of
generative models: the number of possible attribute combinations grows
multiplicatively, but training data is finite.

\subsubsection*{Real-World Story: The Prompt Engineering Economy}

Small changes in wording---``a photo of a cat'' versus ``a professional
photograph of a cat''---can produce dramatically different outputs.  This
sensitivity is not a bug but a direct consequence of how text embeddings
work: nearby words occupy nearby points in embedding space, and even small
shifts in position alter the guidance signal.  \emph{Lesson:} prompt
sensitivity is an architectural property, not a user error.


% ══════════════════════════════════════════════════════════════════════════
\section{The Co-Pilot Reference Card}
\label{sec:reference}

Each session adds one row to a cumulative reference table.  By the end of
Session~5, participants have a complete cheat sheet:

\begin{center}
\renewcommand{\arraystretch}{1.35}
\small
\begin{tabular}{@{} >{\bfseries}p{1.6cm} p{3.2cm} p{3.2cm} p{4cm} @{}}
\toprule
Session & Component & Metaphor & What It Does \\
\midrule
1 & Images \& Text & Car parts \& map & Represent visual and textual information as numbers \\
2 & CLIP & Co-pilot reads map & Aligns image and text representations in a shared space \\
3 & Diffusion & Driver navigates & Generates images by iteratively removing noise \\
4 & Cross-Attention & Driver asks co-pilot & Injects text guidance into the image generation process \\
5 & Full Pipeline & The road trip & Wires all components into a working text-to-image system \\
\bottomrule
\end{tabular}
\end{center}


% ══════════════════════════════════════════════════════════════════════════
\section{Five Questions to Ask an AI Image Generation Vendor}
\label{sec:questions}

\begin{enumerate}
  \item \textbf{Data representation.}  What image resolution, colour
        space, and file format does your model produce---and do these
        match our downstream requirements?
  \item \textbf{Training data.}  How large and diverse is your training
        dataset, and what bias audits have been conducted?
  \item \textbf{Architecture.}  What is the architecture of your
        denoiser, and what are its known failure modes (e.g., text
        rendering, fine detail, attribute binding)?
  \item \textbf{Diagnostics.}  Can you provide attention maps or other
        interpretability tools to diagnose failed generations?
  \item \textbf{Reproducibility.}  Does your system log random seeds and
        model versions for every generated asset, enabling deterministic
        replay?
\end{enumerate}


% ══════════════════════════════════════════════════════════════════════════
\section{Vendor Evaluation: Decision Tree}
\label{sec:decision}

When adopting AI image generation, organisations face three broad
options:

\begin{description}[style=nextline, leftmargin=1.5cm]
  \item[API Service] Use a third-party API (e.g., OpenAI, Stability AI).
    Best when: speed to market is paramount, volume is moderate, and
    customisation needs are low.
  \item[Fine-Tuned Model] Take a pre-trained open model and fine-tune
    it on proprietary data.  Best when: domain-specific quality matters
    (e.g., your brand's visual style) and you have in-house ML capacity.
  \item[Self-Hosted Model] Train or deploy a model entirely on your own
    infrastructure.  Best when: data privacy is non-negotiable, you need
    full control over the model, and you have significant engineering
    resources.
\end{description}

The right choice depends on your organisation's constraints along four
axes: \textbf{data sensitivity}, \textbf{customisation depth},
\textbf{engineering capacity}, and \textbf{cost structure}
(per-generation API fees vs.\ upfront infrastructure investment).


% ══════════════════════════════════════════════════════════════════════════
\section{Jargon Decoder}
\label{sec:jargon}

\begin{center}
\renewcommand{\arraystretch}{1.3}
\small
\begin{tabular}{@{} >{\bfseries}p{3.5cm} p{9cm} @{}}
\toprule
Term & Plain-Language Definition \\
\midrule
Tensor & A multi-dimensional grid of numbers (an image is a 3-D tensor) \\
Embedding & A list of numbers representing a word, sentence, or image in a form computers can process \\
Cosine similarity & A score measuring how similar two embeddings are (1 = identical direction, 0 = unrelated) \\
CLIP & A model that learns to place matching images and texts near each other in a shared number space \\
Contrastive learning & A training strategy that pushes matches together and mismatches apart \\
Diffusion & A generative process that creates images by iteratively removing noise from random static \\
U-Net & A neural network architecture with skip connections, well suited for image denoising \\
Cross-attention & A mechanism allowing one type of data (images) to query another (text) for guidance \\
Guidance scale & A parameter controlling how strongly the text prompt influences image generation \\
Noise schedule & The plan governing how quickly noise is added or removed at each step \\
\bottomrule
\end{tabular}
\end{center}


% ══════════════════════════════════════════════════════════════════════════
\section{Red Flags When Evaluating Vendors}
\label{sec:redflags}

\begin{enumerate}
  \item \textbf{No bias audit.}  The vendor cannot describe what biases
        exist in their training data or what mitigations are in place.
  \item \textbf{No reproducibility.}  Generated outputs cannot be
        recreated deterministically (no seed logging, no version
        tracking).
  \item \textbf{Black-box failures.}  When generation fails, the vendor
        offers no diagnostic tools (e.g., attention maps) and simply
        suggests ``try a different prompt.''
  \item \textbf{Benchmark-only evaluation.}  Quality claims are based
        solely on standard benchmarks (FID, CLIP score) without
        evaluation on your domain-specific data.
  \item \textbf{No content-provenance metadata.}  Generated images carry
        no watermark or C2PA metadata, making it impossible to
        distinguish them from real photographs downstream.
\end{enumerate}


% ══════════════════════════════════════════════════════════════════════════
\section{Conclusion}
\label{sec:conclusion}

The core architecture of text-to-image generation rests on three
learnable components---a text--image alignment model (CLIP), a denoising
network (U-Net with diffusion), and a conditioning mechanism
(cross-attention)---plus a set of carefully chosen hyper-parameters
(noise schedule, guidance scale, seed).

None of these pieces is, on its own, mysterious.  Images are grids of
numbers.  Text is lists of numbers.  CLIP aligns the two.  Diffusion
learns to remove noise.  Cross-attention lets text steer the denoiser.
The compound system produces results that feel like magic, but every step
is auditable, tuneable, and---as you have now seen---buildable from
scratch.

Understanding these building blocks transforms the conversation with
engineering teams and vendors from ``Can the AI do this?''\ to ``Which
component is responsible for this limitation, and what would it take to
fix it?''

\vfill

\begin{center}
\small\color{cGray}
\textit{These lecture notes accompany the interactive course
``Stable Diffusion from Scratch: A Decision-Maker's Guide.''}\\
\textit{All figures and interactive elements are available in the
companion Jupyter notebooks.}
\end{center}

\end{document}
